\documentclass[10pt,a4paper]{amsart}
\usepackage[margin=19mm]{geometry}
\usepackage{amsmath,amssymb,mathtools}
\usepackage[scr=rsfs,cal=euler]{mathalfa}
\usepackage{xfrac}
\usepackage[unicode,hidelinks,bookmarks=false]{hyperref}
\newcommand{\ie}{i.e.\ }
\newcommand{\cf}{cf.\ }
\newcommand{\resp}{respectively\ }
\newcommand{\Subsection}{\subsection}
\providecommand{\nobreakdash}{\nobreak\hbox{-}}
\setlength{\emergencystretch}{2em}
\multlinegap0pt
\usepackage{amssymb}
\usepackage{xcolor}
\newtheorem{theorem}{Theorem}
\newtheorem{proposition}{Proposition}
\newtheorem{lemma}{Lemma}
\newcommand{\R}{\mathbb{R}}
\title{Blow-up and blow-up-time estimates for a singular pseudo-parabolic equation with a space--time variable exponent}
\author{Nguyen Thanh Tung, Le Xuan Truong, Tan Duc Do, Nguyen Ngoc Trong}
\date{Source: arXiv:2506.04498v3 (CC BY 4.0)}
\begin{document}
\maketitle

\noindent\emph{Source and licence.} Blow-up and blow-up-time estimates for a singular pseudo-parabolic equation with a space--time variable exponent. Version arXiv:2506.04498v3 (CC BY 4.0), distributed by arXiv under the Creative Commons Attribution 4.0 International licence, http://creativecommons.org/licenses/by/4.0/, as recorded in the arXiv licence metadata for this exact version and shown on the abstract page https://arxiv.org/abs/2506.04498v3. Full licence notice: \texttt{licenses/arxiv-2506.04498-CC-BY-4.0.txt}.\par\smallskip

\noindent\textit{Editorial scope.} Counted unit: the article's theorem thm:positive (source0.tex lines 220--246). The article's complete original proof of it is delivered with it (source0.tex lines 666--823, one \texttt{proof} environment of 158 lines, which the article places away from the statement). No mathematics is written, repaired or re-derived: the statement and the proof are the article's own lines, in the article's own notation, and no retained line is substituted or re-flowed.  Retained uncounted as the source-local context the proof needs: lines 99--106; lines 114--123; lines 174--176; lines 322--327; lines 491--507; lines 562--574; lines 582--592.  Nothing was omitted from any retained span, so the package carries no omission record.  No retained line contains a citation, so the wrapper carries no bibliography and no bibliography entry.  No retained line contains a figure environment, an image-inclusion command or drawing code, so no image is delivered and the package declares no asset dependency.  Editorial formatting only, in addition: an AMS \texttt{amsart} preamble that restates the article's own declarations and macros used by the retained lines, namely the environments \texttt{theorem}, \texttt{proposition}, \texttt{lemma} on separate counters, together with the standard packages those lines need.  The retained environments print in this document as Theorem 1, Proposition 1, Lemma 1 in order of appearance, whereas the article prints the same items with the numbering of its own full document.  The complete native source of the article as distributed in the arXiv e-print for this exact version is bundled unchanged as \texttt{sources/179346/source0.tex}.
\begin{equation} \label{eq:p-assumption}
	\begin{cases}
		p \in C^1(\overline{Q}), \quad p_t \ge 0,
		\\[2mm]
		\displaystyle
		2 < p^- \le p(\cdot) \le p^+ < 2^*:=\frac{2n}{n-2}.
	\end{cases}
\end{equation}

\begin{equation} \label{eq:k-assumption}
	\begin{cases}
		k \in C^1[0,\infty), 
		\quad k(0) > 0
		\quad \mbox{and} \quad
		k'(t) \ge 0  \,\, \mbox{ for all } t \in [0,\infty),
		\\[2mm]
		k_\infty := \displaystyle\lim_{t\to\infty} k(t) < \infty.
	\end{cases}
\end{equation}

	\begin{equation} \label{eq:blowup}
		\limsup_{t\to(T^*)^-}L(t)=\infty.
	\end{equation}

\begin{proposition}\label{prop:local}
For every $u_0\in W^{1,2}_0(\Omega)$ there exists a unique maximal weak solution $u$ of $(P)$ on an interval $[0,T^*)$, $T^*\in(0,\infty]$. Moreover, if $T^*<\infty$, then
\[
\limsup_{t\to(T^*)^-}L(t)=\infty.
\]
\end{proposition}

\begin{lemma} \label{lem:concavity}
	Let $\theta>0$ and let $\psi$ be positive and continuously differentiable on $[0,T)$, with $\psi'$ locally absolutely continuous, such that
	$\psi(0)>0$, $\psi'(0)>0$, and
	\[
	\psi(t)\psi''(t)-(1+\theta)(\psi'(t))^2\ge0,
	\qquad \text{for a.e. }0<t<T.
	\]
	Then
	\[
	t<\frac{\psi(0)}{\theta\psi'(0)}
	\qquad\text{for every }0<t<T.
	\]
	Consequently,
	\[
	T\le\frac{\psi(0)}{\theta\psi'(0)}.
	\]
\end{lemma}

	\begin{align}
	K'(t)
	={}&-\frac{d}{dt}\int_\Omega\frac{k_\infty}{p(x,t)}\,dx
	-\frac{d}{dt}J(u(t),t) \notag\\
	={}&\left\|\frac{u_t(t)}{|x|}\right\|_2^2
	+\|\nabla u_t(t)\|_2^2
	+k'(t)\int_\Omega\frac{|u(x,t)|^{p(x,t)}}{p(x,t)}\,dx \notag\\
	&\quad +k(t)\mathfrak{P}(t)
	-\frac{d}{dt}\int_\Omega\frac{k_\infty}{p(x,t)}\,dx \notag\\
	\ge{}&\left\|\frac{u_t(t)}{|x|}\right\|_2^2
	+\|\nabla u_t(t)\|_2^2\ge0.
	\label{eq:K-derivative}
	\end{align}

	\begin{align}
	L'(t)
	&=-I(u(t),t) \notag\\
	&=-\|\nabla u(t)\|_2^2
	+k(t)\int_\Omega |u(x,t)|^{p(x,t)}\,dx \notag\\
	&\ge \frac{p^- -2}{2}\|\nabla u(t)\|_2^2-p^-J(u(t),t) \notag\\
	&\ge -p^-J(u(t),t) \notag\\
	&=p^-\left[K(t)+\int_\Omega\frac{k_\infty}{p(x,t)}\,dx\right]
	\ge p^-K(t)>0
	\label{eq:L-prime-negative}
	\end{align}

\begin{theorem} \label{thm:positive}
	Let $n \in \{3,4,5,\ldots\}$ and $\Omega \subset \R^n$ be open bounded with Lipschitz boundary.
	Let $p$, $k$ satisfy \eqref{eq:p-assumption} and \eqref{eq:k-assumption} respectively, and set
	\[
	C_1:=\frac{p^-\,(1+H_n)}{p^- -2},
	\qquad H_n:=\frac{4}{(n-2)^2}.
	\]
	Suppose that the initial modified energy is nonnegative and satisfies
	\[
	0\le C_1E(u_0,0)<L(0),
	\qquad
	L(0)=\frac12\left(
	\left\|\frac{u_0}{|x|}\right\|_{L^2(\Omega)}^2
	+\|\nabla u_0\|_{L^2(\Omega)}^2
	\right).
	\]
	Let $u$ be the maximal weak solution to $(P)$ and define
	\[
	M_0:=L(0)-C_1E(u_0,0)>0.
	\]
	Then $u$ blows up at a finite time $T^*$ satisfying
	\[
	T^*\le
	\frac{8C_1(p^- -1)L(0)}
	{p^-(p^- -2)^2M_0}.
	\]
\end{theorem}

\begin{proof}[Proof of Theorem \ref{thm:positive}]
	Let $T^*\in(0,\infty]$ be the maximal existence time of $u$, and write
	\[
	D(t):=\left\|\frac{u_t(t)}{|x|}\right\|_2^2+\|\nabla u_t(t)\|_2^2.
	\]
	The estimate \eqref{eq:K-derivative}, whose derivation does not use the sign of $E(u_0,0)$, yields
	\begin{equation}\label{eq:K-positive}
	K'(t)\ge D(t)\ge0,
	\qquad K(t):=-E(u(t),t).
	\end{equation}
	In particular,
	\begin{equation}\label{eq:K-integrated}
	K(t)\ge K(0)+\int_0^tD(s)\,ds
	=-E(u_0,0)+\int_0^tD(s)\,ds.
	\end{equation}

	By Hardy's inequality,
	\[
	\left\|\frac{u(t)}{|x|}\right\|_2^2\le H_n\|\nabla u(t)\|_2^2,
	\]
	and hence
	\begin{equation}\label{eq:gradient-controls-L}
	\|\nabla u(t)\|_2^2
	\ge \frac{2}{1+H_n}L(t).
	\end{equation}
	As in \eqref{eq:L-prime-negative}, but retaining the gradient contribution, we have
	\begin{align}
	L'(t)
	&\ge \frac{p^- -2}{2}\|\nabla u(t)\|_2^2-p^-J(u(t),t) \notag\\
	&=\frac{p^- -2}{2}\|\nabla u(t)\|_2^2
	+p^-K(t)+p^-\int_\Omega\frac{k_\infty}{p(x,t)}\,dx \notag\\
	&\ge \frac{p^- -2}{1+H_n}L(t)+p^-K(t). \label{eq:L-prime-positive}
	\end{align}
	Set
	\[
	a:=\frac{p^- -2}{1+H_n},
	\qquad
	C_1=\frac{p^-}{a}=\frac{p^-(1+H_n)}{p^- -2},
	\qquad
	M(t):=L(t)+C_1K(t).
	\]
	Then \eqref{eq:L-prime-positive} gives
	\[
	L'(t)\ge aM(t),
	\]
	while \eqref{eq:K-positive} implies
	\[
	M'(t)=L'(t)+C_1K'(t)\ge L'(t)\ge aM(t).
	\]
	Since
	\[
	M(0)=L(0)-C_1E(u_0,0)=M_0>0,
	\]
	Gronwall's inequality gives $M(t)\ge M_0e^{at}>0$. Therefore
	\begin{equation}\label{eq:L-increasing}
	L'(t)>0,
	\qquad L(t)\ge L(0),
	\qquad 0<t<T^*.
	\end{equation}

	We now apply a corrected concavity argument. Fix $0<\tau<T^*$ and choose
	\begin{equation}\label{eq:beta-choice}
	0<\beta<\frac{p^-M_0}{C_1(p^- -1)}.
	\end{equation}
	For a parameter $\sigma>0$ to be fixed below, define
	\[
	F(h):=\int_0^hL(s)\,ds+(\tau-h)L(0)+\frac{\beta}{2}(h+\sigma)^2,
	\qquad 0\le h\le\tau.
	\]
	Then $F(h)>0$ and, using $L'(s)=\langle u(s),u_t(s)\rangle_{\mathcal H}$ for the natural Hilbert product
	\[
	\langle v,w\rangle_{\mathcal H}
	:=\left(\frac{v}{|x|},\frac{w}{|x|}\right)_{L^2}
	+(\nabla v,\nabla w)_{L^2},
	\]
	we have
	\[
	F'(h)=\int_0^h\langle u(s),u_t(s)\rangle_{\mathcal H}\,ds
	+\beta(h+\sigma),
	\qquad
	F''(h)=L'(h)+\beta.
	\]
	Cauchy--Schwarz in $L^2(0,h;\mathcal H)\times\mathbb R$ gives
	\begin{align}
	(F'(h))^2
	&\le\left(2\int_0^hL(s)\,ds+\beta(h+\sigma)^2\right)
	\left(\int_0^hD(s)\,ds+\beta\right)\notag\\
	&\le 2F(h)\left(\int_0^hD(s)\,ds+\beta\right). \label{eq:F-CS}
	\end{align}

	Next, \eqref{eq:K-integrated} and $-J(u(h),h)=K(h)+\int_\Omega k_\infty/p(x,h)\,dx\ge K(h)$ yield
	\begin{align}
	L'(h)
	&\ge \frac{p^- -2}{2}\|\nabla u(h)\|_2^2+p^-K(h)\notag\\
	&\ge \frac{p^- -2}{2}\|\nabla u(h)\|_2^2
	-p^-E(u_0,0)+p^-\int_0^hD(s)\,ds. \label{eq:L-prime-concavity}
	\end{align}
	Put
	\[
	\theta:=\frac{p^- -2}{2}>0.
	\]
	Since $1+\theta=p^-/2$, inequalities \eqref{eq:F-CS} and \eqref{eq:L-prime-concavity} give
	\begin{align*}
	&F(h)F''(h)-(1+\theta)(F'(h))^2\\
	&\quad\ge F(h)\left[
	F''(h)-p^-\left(\int_0^hD(s)\,ds+\beta\right)
	\right]\\
	&\quad\ge F(h)\left[
	\frac{p^- -2}{2}\|\nabla u(h)\|_2^2
	-p^-E(u_0,0)-(p^- -1)\beta
	\right]\\
	&\quad\ge F(h)\left[
	\frac{p^- -2}{1+H_n}L(0)
	-p^-E(u_0,0)-(p^- -1)\beta
	\right]\\
	&\quad=F(h)\left[
	\frac{p^-}{C_1}M_0-(p^- -1)\beta
	\right]>0,
	\end{align*}
	where we used \eqref{eq:gradient-controls-L}, \eqref{eq:L-increasing}, and \eqref{eq:beta-choice}.

	Lemma \ref{lem:concavity} therefore yields
	\[
	\tau\le\frac{F(0)}{\theta F'(0)}
	=\frac{2\tau L(0)+\beta\sigma^2}
	{(p^- -2)\beta\sigma}.
	\]
	Choose
	\begin{equation}\label{eq:sigma-choice}
	\sigma>\frac{2L(0)}{(p^- -2)\beta}.
	\end{equation}
	Then
	\[
	\tau\le
	\frac{\beta\sigma^2}
	{(p^- -2)\beta\sigma-2L(0)}.
	\]
	The right-hand side is minimized at
	\[
	\sigma=\frac{4L(0)}{(p^- -2)\beta},
	\]
	which is admissible in \eqref{eq:sigma-choice}, and hence
	\begin{equation}\label{eq:tau-beta}
	\tau\le\frac{8L(0)}{(p^- -2)^2\beta}.
	\end{equation}
	Because \eqref{eq:tau-beta} holds for every $\beta$ satisfying \eqref{eq:beta-choice}, letting
	\[
	\beta\uparrow\frac{p^-M_0}{C_1(p^- -1)}
	\]
	gives
	\[
	\tau\le
	\frac{8C_1(p^- -1)L(0)}
	{p^-(p^- -2)^2M_0}.
	\]
	Finally let $\tau\uparrow T^*$. We obtain the asserted finite upper bound for $T^*$.
	Proposition \ref{prop:local} then gives \eqref{eq:blowup}, so the solution blows up at $T^*$.
\end{proof}

\end{document}
